\section{Constructing and certifying approximate corrections}
\label{sec:certification}

The useful residual is the component of acceleration that changes the
distribution.  This component has a canonical definition even when its
orthogonal complement does not have a classical flow.  For a probability
measure $p$, write
\[
 \mathcal G_p=\overline{\{\nabla\phi:\phi\in C_c^\infty(\R^d)\}}^{L^2(p)}
\]
and let $\Pi_p$ be the orthogonal projection onto this space.
Its orthogonal complement consists exactly
of the fields with zero weighted divergence.  This is the standard tangent
space decomposition in Wasserstein geometry~\citep[Chapter~8]{ambrosio2005gradient};
the contribution below is its quantitative use in the Euler error bound.
Define
\begin{equation}
 \varepsilon_t=\|\Pi_{p_t}\acc_t\|_{L^2(p_t)}
 =\sup_{\substack{\phi\in C_c^\infty\\\|\nabla\phi\|_{L^2(p_t)}\le1}}
       \left|\int\acc_t\cdot\nabla\phi\,\dd p_t\right|.
 \label{eq:canonical-residual}
\end{equation}
Thus $\varepsilon_t=0$ is equivalent to the acceleration condition.
No density or Poincar\'e inequality is needed for this definition.

Here canonical refers to the projection for the chosen generator $u$,
not an intrinsic acceleration of the probability path. For the expanding
circular Gaussian, $u(x)=ax$ gives $\Pi_p\acc=a^2x$, whereas
$u(x)=a(I+J)x$ gives $\Pi_p\acc=0$ on the same path. Intrinsic
Wasserstein covariant acceleration instead first uses the minimum-energy
gradient velocity $\Pi_pu$.

\subsection{A residual theorem without a regular projected flow}

\begin{theorem}[Projected acceleration controls Euler error]
\label{thm:projected-residual}
Keep the smoothness, spatial Lipschitz, growth, and moment assumptions of
Theorem~\ref{thm:distribution_order}, including the common Lipschitz constant
$L_a$ of $\acc_t$, but do not impose the acceleration condition.  Set
\begin{align}
 E_H&=\exp(L_aH^2/2),\qquad B_H=(1+LH)E_H,\nonumber\\
 D_H&=\frac{C_j}{6}+E_H\left(\frac{LC_a}{2}
                         +\frac{HL_aC_a}{8}\right).
 \label{eq:projected-constants}
\end{align}
Then on every finite increasing grid,
\begin{align}
 \max_k W_2(\widehat p_k,p_{t_k})
 &\le e^{LT}\left(D_H\sum_jh_j^3
                   +\frac{B_H}{2}\sum_jh_j^2\varepsilon_{t_j}\right)
 \label{eq:projected-bound}\\
 &\le e^{LT}T\left(D_HH^2+\frac{B_H}{2}\varepsilon H\right)
 .\nonumber
\end{align}
The second line assumes $\sup_t\varepsilon_t\le\varepsilon$.
In particular, the projected solenoidal field
$\acc_t-\Pi_{p_t}\acc_t$ need not be continuous or Lipschitz.
\end{theorem}

Proposition~\ref{prop:robustness} required a regular flow for the
distribution-preserving part of acceleration.  Theorem~\ref{thm:projected-residual}
instead follows the full acceleration field and bounds the change of its
distribution by the projected residual.  The price is the explicit factor
$B_H=1+O(H)$ and a different cubic constant; no elliptic regularity of the
projection is required.  Appendix~\ref{app:certification-proofs} proves the
result.  A certificate $\varepsilon_{t_j}\le\kappa H$ at every Euler node
gives an $O(H^2)$ bound when the other constants are uniform.  The certificate
may belong to one fixed field or to a family of fields indexed by $H$; these
are different assertions.  A small value only at a training grid does not
control a different evaluation grid without an interpolation bound.

\begin{corollary}[Exact local Wasserstein coefficient]
\label{cor:sharp-local-residual}
Under Theorem~\ref{thm:projected-residual}, at every fixed $t<T$,
\begin{equation}
 \lim_{h\downarrow0}\frac{W_2((E_{t,h})_\#p_t,p_{t+h})}{h^2}
 =\tfrac12\|\Pi_{p_t}\acc_t\|_{L^2(p_t)}.
 \label{eq:sharp-local-residual}
\end{equation}
\end{corollary}
The projected residual is therefore the exact leading local error
coefficient.  This strengthens the local necessity result, but does not
rule out cancellation between different steps at the final time.

\subsection{Upper and lower certificates for the residual}

A finite test space detects only part of the gradient projection.  Its norm
is a lower bound, so reducing it alone cannot certify a small residual.
The following identity supplies both sides using a gradient approximation
and an independently admissible solenoidal field.

\begin{proposition}[Primal and dual residual certificates]
\label{prop:residual-certificates}
Fix $p$ and $a\in L^2(p;\R^d)$.  For any $q\in\mathcal G_p$ and
$w\in\mathcal G_p^\perp$, put
\begin{align*}
 \mathcal L&=2\langle a,q\rangle_{L^2(p)}-\|q\|_{L^2(p)}^2,\\
 \mathcal U&=\|a-w\|_{L^2(p)}^2.
\end{align*}
Then
\begin{align}
 \max(0,\mathcal L)&\le\|\Pi_pa\|_{L^2(p)}^2\le\mathcal U,\nonumber\\
 \mathcal U-\mathcal L&=\|a-w-q\|_{L^2(p)}^2.
 \label{eq:primal-dual-certificate}
\end{align}
For an exact Galerkin solution
$q_N=\sum_i c_i\nabla\phi_i$, where
$A_{ij}=\langle\nabla\phi_i,\nabla\phi_j\rangle$ and
$b_i=\langle a,\nabla\phi_i\rangle$, $Ac=b$ gives
$\mathcal L=\|q_N\|^2=b^\top A^{-1}b$ when $A$ is positive definite.
The general certificate remains valid for inexact coefficients.
\end{proposition}

This Hilbert-space identity turns a numerical projection into an upper
certificate for Eq.~\eqref{eq:projected-bound}.  For fixed acceleration,
gradient and solenoidal linear bases give two quadratic least-squares
problems; the gap checks their resolution.  An explicit admissible $w$
already gives an upper bound, so a Poisson solve is optional.

A weighted Poisson formulation, an a posteriori Poincar\'e estimate, and
explicit bounds for quadrature and acceleration approximation errors appear
in Appendix~\ref{app:certification-proofs}.  A raw divergence penalty or
Monte Carlo standard error alone is not a deterministic residual certificate.

\subsection{A finite-dimensional construction with explicit scope}

In two dimensions let $J$ denote the quarter-turn matrix.  For a known
positive density, stream potentials give the established gauge
parameterization~\citep{horvat2024gauge}
\begin{align}
 r_\theta(t,x)&=p_t(x)^{-1}J\nabla_x\psi_\theta(t,x),\nonumber\\
 \psi_\theta&=\sum_{\ell=1}^m\theta_\ell\psi_\ell,
 \qquad \nabla\cdot(p_tr_\theta)=0.
 \label{eq:compact-stream-gauge}
\end{align}
Compactly supported potentials avoid dividing by arbitrarily small density
in the tails.  The proof appendix gives the antisymmetric-tensor version in higher
dimensions. A finite family need not represent every admissible field.

For a base field satisfying the theorem's bounds, a positive $C^4$ density,
uniformly compactly supported $C^4$ potentials and a compact parameter set,
the corrected fields preserve the
exact path and have uniform regularity bounds.  Their objective
$\mathcal E(\theta)=\sup_t\|\Pi_{p_t}\acc_{u_\theta,t}\|$ is
Lipschitz and attains a minimum.  Proposition~\ref{prop:compact-gauge-search}
in the proof appendix gives the assumptions and a certified finite-net
search guarantee: parameter mesh $\delta$ and objective certificate width
$\eta$ achieve residual at most
$\min_\Theta\mathcal E+K_\Theta\delta+\eta$.


This finite-dimensional guarantee does not imply efficiency or convergence
of a local nonlinear optimizer.  A positive minimum in one family does
not prove path infeasibility.  The appendix also certifies the continuous
time supremum from a time grid.

\subsection{Separating path error from integration error}

Let $p_t$ now be a target path generated by $v$, and let $u$ be the field
actually integrated.  Its exact path need not equal $p_t$.  Besides the
acceleration residual under the target, define the path defect
\begin{equation}
 \beta_t=\|\Pi_{p_t}(u_t-v_t)\|_{L^2(p_t)}.
 \label{eq:path-defect}
\end{equation}
This defect ignores velocity differences that preserve the target
continuity equation.  Both diagnostics in the next theorem are evaluated under the target
$p_t$; no density substitution for an unknown learned path is made.

\begin{theorem}[Target-path, residual, and integration errors]
\label{thm:target-path-bound}
Let $p_t$ be generated by a smooth, uniformly spatially Lipschitz $v$ from
$p_0\in\mathcal P_2(\R^d)$, with
$\sup_t\|v_t\|_{L^2(p_t)}<\infty$.  Let $u$ satisfy the smoothness,
growth, and Lipschitz assumptions of Theorem~\ref{thm:projected-residual}.
Write $F^u_{t,s}$ for its exact flow and assume
\begin{align*}
 C_a&=\sup_t\|\acc_{u,t}\|_{L^2(p_t)}<\infty,\\
 C_j^\diamond&=
 \sup_{0\le t\le s\le T}
 \|\jerk_u(s,F^u_{t,s}(\cdot))\|_{L^2(p_t)}<\infty.
\end{align*}
Suppose $\beta_t$ is integrable.  Let $\widehat q_k$ be Euler's laws for
$u$ from any $q_0\in\mathcal P_2$.  With
$\varepsilon_t=\|\Pi_{p_t}\acc_{u,t}\|$ and with $D_H$ in
Eq.~\eqref{eq:projected-constants} using $C_j^\diamond$ in place of
$C_j$,
\begin{align}
 \max_k W_2(\widehat q_k,p_{t_k})
 \le e^{LT}\bigg[&W_2(q_0,p_0)+D_H\sum_jh_j^3\nonumber\\
                 &+\frac{B_H}{2}\sum_jh_j^2\varepsilon_{t_j}
                  +\int_0^T\beta_t\,\dd t\bigg].
 \label{eq:target-path-bound}
\end{align}
\end{theorem}

The mixed jerk assumption concerns exact $u$-trajectories started from
the target at any time.  It is automatic from a uniform linear-growth
bound on $\jerk_u$ and bounded target second moments, but does not follow
from a training loss alone.  If $u$ generates $p_t$, then
$C_j^\diamond=C_j$ and $\beta_t=0$, recovering
Theorem~\ref{thm:projected-residual}.

For conditional flow matching with $v_t(x)=\E[Y_t\mid X_t=x]$ and
$\Law(X_t)=p_t$, the population excess squared regression loss is
$\mathcal R_{\rm excess}=\int_0^T\|u_t-v_t\|_{L^2(p_t)}^2\,\dd t$.
Projection contraction gives
$\int_0^T\beta_t\,\dd t\le\sqrt{T\mathcal R_{\rm excess}}$.
These population and exact-flow facts are standard~\citep{lipman2023flow,benton2024error};
Eq.~\eqref{eq:target-path-bound} combines them with the Euler residual.
There is no finite-sample learning guarantee.  An exactly target-solenoidal
correction leaves $\beta_t$ unchanged, but can change the other constants.

For a stream correction formed with an approximate density, the resulting
path defect can be bounded using the score discrepancy and a Poincar\'e
constant; Eq.~\eqref{eq:score-path-defect} gives the precise bound.
